\documentclass[10pt,a4paper]{amsart}
\usepackage[margin=19mm]{geometry}
\usepackage{amsmath,amssymb,mathtools}
\usepackage[scr=rsfs,cal=euler]{mathalfa}
\usepackage{xfrac}
\usepackage[unicode,hidelinks,bookmarks=false]{hyperref}
\newcommand{\ie}{i.e.\ }
\newcommand{\cf}{cf.\ }
\newcommand{\resp}{respectively\ }
\newcommand{\Subsection}{\subsection}
\providecommand{\nobreakdash}{\nobreak\hbox{-}}
\setlength{\emergencystretch}{2em}
\multlinegap0pt
\usepackage{amssymb}
\usepackage{mathtools}
\usepackage{bm}
\usepackage{tikz}
\usepackage{xcolor}
\usepackage[normalem]{ulem}
\newtheorem{lemma}{Lemma}
\newcommand{\RePart}{\operatorname{Re}}
\newcommand{\T}{\mathbb T}
\newcommand{\tr}{\operatorname{tr}}
\title{A Complete Classification of Complex Hadamard Matrices of Order Six}
\author{Mateo C\'ardenes Wuttig, Joseph Tindall}
\date{Source: arXiv:2608.18053v2 (CC BY 4.0)}
\begin{document}
\maketitle

\noindent\emph{Source and licence.} A Complete Classification of Complex Hadamard Matrices of Order Six. Version arXiv:2608.18053v2 (CC BY 4.0), distributed by arXiv under the Creative Commons Attribution 4.0 International licence, http://creativecommons.org/licenses/by/4.0/, as recorded in the arXiv licence metadata for this exact version and shown on the abstract page https://arxiv.org/abs/2608.18053v2. Full licence notice: \texttt{licenses/arxiv-2608.18053-CC-BY-4.0.txt}.\par\smallskip

\noindent\textit{Editorial scope.} Counted unit: the article's lemma lem:rowcolumn (source0.tex lines 1535--1542). The article's complete original proof of it is delivered with it (source0.tex lines 1544--1568, one \texttt{proof} environment of 25 lines). No mathematics is written, repaired or re-derived: the statement and the proof are the article's own lines, in the article's own notation, and no retained line is substituted or re-flowed.  No further source-local context is needed: every cross-reference of every retained line resolves inside the two delivered spans.  Nothing was omitted from any retained span, so the package carries no omission record.  No retained line contains a citation, so the wrapper carries no bibliography and no bibliography entry.  No retained line contains a figure environment, an image-inclusion command or drawing code, so no image is delivered and the package declares no asset dependency.  Editorial formatting only, in addition: an AMS \texttt{amsart} preamble that restates the article's own declarations and macros used by the retained lines, namely the environments \texttt{lemma} on separate counters, together with the standard packages those lines need.  The retained environments print in this document as Lemma 1 in order of appearance, whereas the article prints the same items with the numbering of its own full document.  The complete native source of the article as distributed in the arXiv e-print for this exact version is bundled unchanged as \texttt{sources/207963/source0.tex}.
\begin{lemma}[Row--column invariant identity]
\label{lem:rowcolumn}
For every $X\in\T^{3\times3}$,
\begin{equation}
\RePart\tau_{\rm r}(X)=\RePart\tau_{\rm c}(X).
\label{eq:rowcolidentity}
\end{equation}
\end{lemma}

\begin{proof}
Set $A=XX^{\dagger}$ and $C=X^{\dagger}X$.  Both are Hermitian, both have diagonal $(3,3,3)$, and they have the same eigenvalues because $X$ is square. For a Hermitian matrix $M$ with this diagonal, define
\begin{equation}
\sigma(M)=\sum_{1\leq i<j\leq3}|M_{ij}|^2.
\label{eq:sigmam}
\end{equation}
Direct multiplication gives
\begin{equation}
\tr(M^2)=27+2\sigma(M).
\label{eq:tracem2}
\end{equation}
Equality of the spectra of $A$ and $C$ implies equality of their traces of squares, and hence
\begin{equation}
\sigma(A)=\sigma(C).
\label{eq:sigmaequal}
\end{equation}

Expanding a $3\times3$ determinant gives
\begin{equation}
\det M=27-3\sigma(M)
+2\RePart\!\left(M_{12}M_{23}M_{31}\right).
\label{eq:detcubic}
\end{equation}
The equal spectra also give $\det A=\det C$.  Subtracting the two instances of Eq.~\eqref{eq:detcubic} and using Eq.~\eqref{eq:sigmaequal} proves Eq.~\eqref{eq:rowcolidentity}.
\end{proof}

\end{document}
